\section{E-commerce Sales Business}
\label{sec:2-1}

The business sells products through an e-commerce transaction system, and sales generate revenue from customer purchases. At a high level:
\begin{equation*}
    \text{Profit} = \text{Sales Revenue} - \text{Business Costs}
\end{equation*}
Costs may include product/operating costs, marketing costs and retention incentive costs.

\section{Revenue from New and Existing Customers}
\label{sec:2-2}

Sales revenue can be generated by both new customers making their first purchases and existing customers making repeat purchases. New customers contribute to business growth by expanding the customer base, while existing customers may generate additional revenue through continued purchasing activity.

\section{Customer Churn as a Business Problem}
\label{sec:2-3}

When existing customers gradually reduce their purchasing activity or stop purchasing altogether, businesses may lose potential future sales. Understanding the behavior of these customers can help businesses identify opportunities to encourage repeat purchases.

\section{Retention and Marketing Cost Constraints}
\label{sec:2-4}

Retention activities require resources. The company cannot necessarily apply the same retention action to every existing customer because marketing budget, campaign capacity and incentives are limited. Retention resources must be allocated selectively.

\section{Business Objectives}
\label{sec:2-5}

The primary business objective is to support sustainable sales revenue by identifying existing customers who are at risk of churn and enabling more informed customer retention decisions. The specific objectives are:

\begin{itemize}
    \item \textbf{Identify at-risk customers:} recognize existing customers who are likely to stop purchasing within a predefined period.
    \item \textbf{Understand purchasing behavior:} investigate customer characteristics, such as purchase frequency, monetary value and time since the last purchase, that may be associated with churn risk.
    \item \textbf{Support customer prioritization:} provide insights that help the business prioritize customers for potential retention activities.
\end{itemize}

\section{Success Criteria}
\label{sec:2-6}

Success is defined \textbf{relative to simple alternatives} rather than by arbitrary absolute thresholds:

\begin{center}
\begin{tabular}{@{}p{3cm}p{11.5cm}@{}}
    \toprule
    Level & Criterion (evaluated on the unseen test period) \\
    \midrule
    Predictive & Final model PR-AUC exceeds the test-period churn rate (random reference) and the recency-rule baseline. \\
    Targeting & Final model Lift@K and Recall@K at the chosen operating $K$ (e.g.\ 20--30\%) exceed the recency-rule baseline, averaged over weekly scoring dates. \\
    Stability & Top-K performance does not collapse on individual test weeks (report mean $\pm$ standard deviation). \\
    Business scenario & Under the base scenario, estimated net impact at the chosen $K$ is positive and higher than for recency-rule and random targeting at the same $K$. \\
    \bottomrule
\end{tabular}
\end{center}

\section{From Business Objective to Analytical Problem}
\label{sec:2-7}

Increasing revenue requires more repeat-purchase revenue from existing customers and efficient use of retention resources. The analytical direction is therefore to identify existing customers at risk of not purchasing again, estimate churn probability, and prioritize customers for retention activities.

\section{Customer Retention Strategy}
\label{sec:2-8}

\begin{center}
\begin{tabular}{@{}c@{}}
    \textbf{Eligible Existing Customers} \\
    $\downarrow$ \\
    \textbf{Churn Prediction} -- estimate churn probability \\
    $\downarrow$ \\
    \textbf{Risk Ranking} -- rank customers by predicted churn risk \\
    $\downarrow$ \\
    \textbf{Business Rules and Available Capacity} \\
    $\downarrow$ \\
    \textbf{Target Customers for Retention} -- prioritize customers for potential campaigns \\
\end{tabular}
\end{center}

\section{Implementation Roadmap}
\label{sec:2-9}

The approach is implemented as a sequence of steps, each covered by one chapter:

\begin{itemize}
    \item \textbf{Understand the data} (Chapter~\ref{ch:data-understanding}) -- sales and customer behavior, purchase intervals and revenue by purchase cadence, which inform the prediction horizon $H$ and the activity threshold $R_{\max}$.
    \item \textbf{Prepare customer snapshots} (Chapter~\ref{ch:data-preparation}) -- turn transactions into Customer $\times$ Snapshot Date rows with point-in-time features, churn labels and an embargoed temporal split.
    \item \textbf{Model churn risk} (Chapter~\ref{ch:modeling}) -- compare baselines and algorithms and select a final model.
    \item \textbf{Evaluate} (Chapter~\ref{ch:evaluation}) -- predictive, ranking and scenario-based business-impact evaluation on the unseen test period.
    \item \textbf{Deploy} (Chapter~\ref{ch:deployment}) -- weekly batch scoring and Top-K targeting, as a proposed operational design.
\end{itemize}
